\chapter{Differential Equations}
\label{chp:differential_equations}

A differential equation is an equation that relates one or more quantities and
their derivatives. These can often be used to model real world systems.

\begin{note}
In this chapter you will cover:
\begin{itemize}
\item 

Creating a symbolic function;

\item 

Writing a differential equation;

\item 

Solving a differential equation.

\end{itemize}
\end{note}





\section{Tutorial}

You will solve the following problem using a computer to do some of the more
tedious calculations.

A container has volume \(V\) of liquid which is poured in at a rate proportional
to \(e^{-t}\) (where \(t\) is some measurement of time). Initially the container is empty and
after \(t=3\) time units the rate at which the liquid is poured is 15.
\begin{enumerate}

\item 

Show that \(V(t)=\frac{-15e^{3}}{1-e^{3}}(1 - e^{-t})\).

\item 

Obtain the limit \(\lim_{t\to \infty}V(t)\).

\end{enumerate}



You first need to create the differential equation described in the question:

\begin{pyin}
import sympy as sym

t = sym.Symbol("t")
k = sym.Symbol("k")
V = sym.Function("V")

differential_equation = sym.Eq(lhs=sym.diff(V(t), t), rhs=k * sym.exp(-t))
differential_equation
\end{pyin}




\begin{equation*}
\begin{split}\displaystyle \frac{d}{d t} V{\left(t \right)} = k e^{- t}\end{split}
\end{equation*}




In order to solve the differential equation you can write:




\begin{pyin}
sym.dsolve(differential_equation, V(t))
\end{pyin}




\begin{equation*}
\begin{split}\displaystyle V{\left(t \right)} = C_{1} - k e^{- t}\end{split}
\end{equation*}




Note that the question gives an initial condition: ``initially the container
is empty” which corresponds to \(V(0)=0\).
You can pass this to the call to solve the differential equation:




\begin{pyin}
condition = {V(0): 0}
particular_solution = sym.dsolve(differential_equation, V(t), ics=condition)
sym.simplify(particular_solution)
\end{pyin}




\begin{equation*}
\begin{split}\displaystyle V{\left(t \right)} = k - k e^{- t}\end{split}
\end{equation*}




You also know that \(V(3)=15\) which corresponds to the following equation:




\begin{pyin}
equation = sym.Eq(particular_solution.rhs.subs({t: 3}), 15)
equation
\end{pyin}




\begin{equation*}
\begin{split}\displaystyle - \frac{k}{e^{3}} + k = 15\end{split}
\end{equation*}




You can solve this equation to find a value for \(k\):




\begin{pyin}
sym.simplify(sym.solveset(equation, k))
\end{pyin}




\begin{equation*}
\begin{split}\displaystyle \left\{- \frac{15 e^{3}}{1 - e^{3}}\right\}\end{split}
\end{equation*}




which is the required value.


You can use the complete expression for \(V(t)\) to take the limit:




\begin{pyin}
limit = sym.limit((-15 * sym.exp(3) / (1- sym.exp(3))) *  (1 - sym.exp(-t)), t, sym.oo)
limit
\end{pyin}




\begin{equation*}
\begin{split}\displaystyle - \frac{15 e^{3}}{1 - e^{3}}\end{split}
\end{equation*}




This is approximately:




\begin{pyin}
float(limit)
\end{pyin}





\begin{raw}
15.78593544736884
\end{raw}





\begin{note}
In this tutorial you have
\begin{itemize}
\item 

Created a differential equation.

\item 

Obtained the general solution of a differential equation.

\item 

Obtained the particular solution of a differential equation.

\end{itemize}
\end{note}





\section{How to}

\subsection{Create a symbolic function}
\label{\detokenize{tools-for-mathematics/09-differential-equations/how/main:how-to-create-a-symbolic-function}}

To create a symbolic function use \texttt{sympy.Function}.


\begin{api}
sympy.Function("y")
\end{api}



For example:




\begin{pyin}
import sympy as sym

y = sym.Function("y")
y
\end{pyin}





\begin{raw}
y
\end{raw}





You can pass symbolic variables to this symbolic function:




\begin{pyin}
x = sym.Symbol("x")
y(x)
\end{pyin}




\begin{equation*}
\begin{split}\displaystyle y{\left(x \right)}\end{split}
\end{equation*}




You can create the derivative of a symbolic function:




\begin{pyin}
sym.diff(y(x), x)
\end{pyin}




\begin{equation*}
\begin{split}\displaystyle \frac{d}{d x} y{\left(x \right)}\end{split}
\end{equation*}




\subsection{How to create a differential equation}
\label{\detokenize{tools-for-mathematics/09-differential-equations/how/main:how-to-create-a-differential-equation}}

To create a differential equation use \texttt{sympy.Eq}.


\begin{api}
import sympy as sym

y = sym.Function("y")
x = sym.Symbol("x")

equation = sym.Eq(lhs, rhs)
\end{api}



Where \texttt{lhs} and \texttt{rhs} are expressions in \(y\), \(\frac{dy}{dx}\) and \(x\).


For example to create the differential equation: \(\frac{dy}{dx} = \cos(x) y\)
write:




\begin{pyin}
import sympy as sym

y = sym.Function("y")
x = sym.Symbol("x")

lhs = sym.diff(y(x), x)
rhs = sym.cos(x) * y(x)
differential_equation = sym.Eq(lhs, rhs)
differential_equation
\end{pyin}




\begin{equation*}
\begin{split}\displaystyle \frac{d}{d x} y{\left(x \right)} = y{\left(x \right)} \cos{\left(x \right)}\end{split}
\end{equation*}




\subsection{Obtain the general solution of a differential equation}
\label{\detokenize{tools-for-mathematics/09-differential-equations/how/main:how-to-obtain-the-general-solution-of-a-differential-equation}}

To obtain the generation solution to a differential equation use:
\texttt{sympy.dsolve}.


\begin{api}
import sympy as sym

y = sym.Function("y")
x = sym.Symbol("x")

equation = sym.Eq(lhs, rhs)
sym.dsolve(equation, y(x))
\end{api}



For example to solve the differential equation: \(\frac{dy}{dx} = \cos(x) y\) 
write:


\begin{pyin}
import sympy as sym

y = sym.Function("y")
x = sym.Symbol("x")

lhs = sym.diff(y(x), x)
rhs = sym.cos(x) * y(x)
differential_equation = sym.Eq(lhs, rhs)
sym.dsolve(differential_equation, y(x))
\end{pyin}




\begin{equation*}
\begin{split}\displaystyle y{\left(x \right)} = C_{1} e^{\sin{\left(x \right)}}\end{split}
\end{equation*}




\subsection{Obtain the particular solution of a differential equation}
\label{\detokenize{tools-for-mathematics/09-differential-equations/how/main:how-to-obtain-the-particular-solution-of-a-differential-equation}}

To obtain the particular solution to a differential equation use:
\texttt{sympy.dsolve} and pass the initial conditions: \texttt{ics}.


\begin{api}
import sympy as sym

y = sym.Function("y")
x = sym.Symbol("x")

equation = sym.Eq(lhs, rhs)
sym.dsolve(equation, y(x), ics={y(x_0): value})
\end{api}



For example to solve the differential equation: \(\frac{dy}{dx} = \cos(x) y\)
with the condition \(y(5)= \pi\) 
write:




\begin{pyin}
import sympy as sym

y = sym.Function("y")
x = sym.Symbol("x")

lhs = sym.diff(y(x), x)
rhs = sym.cos(x) * y(x)
differential_equation = sym.Eq(lhs, rhs)

condition = {y(5): sym.pi}
sym.dsolve(differential_equation, y(x), ics=condition)
\end{pyin}




\begin{equation*}
\begin{split}\displaystyle y{\left(x \right)} = \pi e^{- \sin{\left(5 \right)}} e^{\sin{\left(x \right)}}\end{split}
\end{equation*}



\begin{note}
The syntax used here is similar to substituting values in to algebraic
expressions (see Section~\ref{sec:substite_a_value_in_to_an_expression}).
\end{note}





\section{Exercises}
\label{\detokenize{tools-for-mathematics/09-differential-equations/exercises/main:exercises}}\label{\detokenize{tools-for-mathematics/09-differential-equations/exercises/main::doc}}
\begin{enumerate}

\item 

Create the following differential equations:
\begin{enumerate}

\item 

\(\frac{dy}{dx} = \cos(x)\)

\item 

\(\frac{dy}{dx} = 1 - y\)

\item 

\(\frac{dy}{dx} = \frac{x - 50}{10}\)

\item 

\(\frac{dy}{dx} = y ^2 \ln (x)\)

\item 

\(\frac{dy}{dx} = (1 + y) ^ 2\)

\end{enumerate}

\item 

Obtain the general solution for the equations in question 1.

\item 

Obtain the particular solution for the equations in question 1 with
the following particular conditions:
\begin{enumerate}

\item 

\(y(0) = \pi\)

\item 

\(y(2) = 3\)

\item 

\(y(50) = 1\)

\item 

\(y(e) = 1\)

\item 

\(y(-1) = 3\)

\end{enumerate}

\item 

The rate of increase of a population (\(p\)) is equal to 1\% of the size of the
population.
\begin{enumerate}

\item 

Define the differential equation that models this situation.

\item 

Given that \(p(0)=5000\) find the population after 5 time units.

\end{enumerate}

\item 

The rate of change of the temperature of a hot drink is proportional to the
difference between the temperature of the drink (\(T\)) and the room temperature (\(T_R\)).
\begin{enumerate}

\item 

Define the differential equation that models this situation.

\item 

Solve the differentia equation.

\item 

Given that \(T(0) = 100\) and the room temperature is \$\(T_R=20\)\$ obtain the
particular solution.

\item 

Use the particular solution to identify how on it will take for the drink
to be ready for consumption (a temperature of 80) given that after 3 time
units \(T(3)=90\).

\end{enumerate}

\end{enumerate}

\section{Further information}
\label{\detokenize{tools-for-mathematics/09-differential-equations/why/main:further-information}}\label{\detokenize{tools-for-mathematics/09-differential-equations/why/main::doc}}

\subsection{How to solve a system of differential equations?}
\label{\detokenize{tools-for-mathematics/09-differential-equations/why/main:how-to-solve-a-system-of-differential-equations}}

Given a system of differential equations like the following:
\begin{equation*}
\begin{split}
    \begin{cases}
        \frac{dx}{dt} =& x - y\\
        \frac{dy}{dt} =& x + y\\
        y(0) =& 250\\
        y(1) =& 300
    \end{cases}
\end{split}
\end{equation*}

You can solve it using \texttt{sym.dsolve} but instead of passing a single differential
equation, pass an iterable of multiple equations:




\begin{pyin}
import sympy as sym


y = sym.Function("y")
x = sym.Function("x")

t = sym.Symbol("t")
alpha = sym.Symbol("alpha")
beta = sym.Symbol("beta")

system_of_equations = (
    sym.Eq(sym.diff(y(t), t), alpha * x(t)),
    sym.Eq(sym.diff(x(t), t), beta * y(t)),
)
conditions = {y(0): 250, y(1): 300}

y_solution, x_solution = sym.dsolve(system_of_equations, ics=conditions, set=True)
x_solution
\end{pyin}




\begin{equation*}
\begin{split}\displaystyle x{\left(t \right)} = - \frac{50 \beta \left(5 e^{\sqrt{\alpha \beta}} - 6\right) e^{\sqrt{\alpha \beta}} e^{- t \sqrt{\alpha \beta}}}{\sqrt{\alpha \beta} \left(e^{2 \sqrt{\alpha \beta}} - 1\right)} + \frac{50 \beta \left(6 e^{\sqrt{\alpha \beta}} - 5\right) e^{t \sqrt{\alpha \beta}}}{\sqrt{\alpha \beta} \left(e^{2 \sqrt{\alpha \beta}} - 1\right)}\end{split}
\end{equation*}






\begin{pyin}
y_solution
\end{pyin}




\begin{equation*}
\begin{split}\displaystyle y{\left(t \right)} = \frac{50 \cdot \left(5 e^{\sqrt{\alpha \beta}} - 6\right) e^{\sqrt{\alpha \beta}} e^{- t \sqrt{\alpha \beta}}}{e^{2 \sqrt{\alpha \beta}} - 1} + \frac{50 \cdot \left(6 e^{\sqrt{\alpha \beta}} - 5\right) e^{t \sqrt{\alpha \beta}}}{e^{2 \sqrt{\alpha \beta}} - 1}\end{split}
\end{equation*}




\subsection{How to solve differential equations numerically?}
\label{\detokenize{tools-for-mathematics/09-differential-equations/why/main:how-to-solve-differential-equations-numerically}}

Some differential equations do not have a closed form solution in terms of
elementary functions. For example, the Airy or \index{Stokes equation}:
\begin{equation*}
\begin{split}
\frac{d^2y}{dx^2} = xy
\end{split}
\end{equation*}

Attempting to solve this with Sympy gives:




\begin{pyin}
import sympy as sym

y = sym.Function("y")
x = sym.Symbol("x")

equation = sym.Eq(sym.diff(y(x), x, 2), x * y(x))
sym.dsolve(equation, y(x))
\end{pyin}




\begin{equation*}
\begin{split}\displaystyle y{\left(x \right)} = C_{1} Ai\left(x\right) + C_{2} Bi\left(x\right)\end{split}
\end{equation*}




which is a linear combination of \(A_i\) and \(B_i\) which are special functions
called the \index{Airy functions} of the first and second kind.


Using \texttt{scipy.integrate} it is possible to solve this differential equation numerically.


First, define a new variable \(u=\frac{dy}{dx}\) so that the second order
differential equation can be expressed as a system of single order differential
equations:
\begin{equation*}
\begin{split}
    \begin{cases}
        \frac{du}{dx}=&xy\\
        \frac{dy}{dx}=&u
    \end{cases}
\end{split}
\end{equation*}

Now define a python function that returns the right hand side of that
system of equations:




\begin{pyin}
def diff(state, x):
    """
    Returns the value of the derivates for a given set of state values (u, y).
    """
    u, y = state
    return x * y, u
\end{pyin}





You can pass this to \texttt{scipy.integrate.odeint} which is a tool that carries out
numerical integration of differential equations. Note, that it is incapable of
dealing with symbolic variables, thus an initial numeric value of \((u, y)\) is
required.




\begin{pyin}
import numpy as np
import scipy.integrate

condition = (.1, -.5)

xs = np.linspace(0, 1, 50)
states = scipy.integrate.odeint(diff, y0=condition, t=xs)

\end{pyin}






\begin{note}
Here, you make use of
\texttt{numpy} to create a collection
of \texttt{x} values over which to carry out the numerical integration.
\end{note}



This returns an array of values of \texttt{states} corresponding to \((u, y)\).





\begin{pyin}
states
\end{pyin}





\begin{raw}
array([[ 0.1       , -0.5       ],
       [ 0.09989617, -0.49795991],
       [ 0.09958578, -0.49592403],
       [ 0.09907053, -0.49389658],
       ...
       [-0.09525243, -0.46835567],
       [-0.10414704, -0.47038996],
       [-0.11327831, -0.47260818],
       [-0.12265169, -0.47501521],
       [-0.13227299, -0.47761605]])
\end{raw}






Figure~\ref{fig:odeint_comparison_plot} shows a plot of the above with a comparison to the exact expected values (obtained
using the \index{Airy functions} of the first and second kind).


\begin{figure}[!htbp]
\begin{center}
\includegraphics[width=.7\textwidth]{./assets/odeint_comparison_plot/main.pdf}
    \caption{Numerical and exact solution to the Stokes differential equation.}
    \label{fig:odeint_comparison_plot}
\end{center}
\end{figure}
